\chapter{Практика 3. Перемежение}

\section{Задание}
Реализовать операцию прямого и обратного перемежения закодированного битового сообщения.

\section{Теоретическое описание}
Перемежение (interleaving) — это процесс перестановки битов в кодовом слове по определенному закону. Основная цель перемежения — борьба с пакетными ошибками. При возникновении пакетной ошибки в канале после обратного перемежения ошибки распределяются по всему кодовому слову, что позволяет эффективно исправлять их с помощью помехоустойчивого кода.

\section{Реализация}
Перемежитель использует псевдослучайную перестановку бит на основе заданного seed. Для обратного перемежения используется сохраненный закон перестановки.

\begin{lstlisting}[language=Python, caption=Перемежитель и деперемежитель]
class Interleaver:
    def __init__(self, seed=42):
        self.seed = seed
        self.permutation = None

    def interleave(self, bits):
        if not bits:
            return bits
        np.random.seed(self.seed)
        n = len(bits)
        self.permutation = np.random.permutation(n)
        interleaved = ['0'] * n
        for i, pos in enumerate(self.permutation):
            interleaved[pos] = bits[i]
        return ''.join(interleaved)

class Deinterleaver:
    def __init__(self, interleaver):
        self.permutation = interleaver.permutation

    def deinterleave(self, bits):
        if not bits or self.permutation is None:
            return bits
        n = len(bits)
        deinterleaved = ['0'] * n
        for i, pos in enumerate(self.permutation):
            deinterleaved[i] = bits[pos]
        return ''.join(deinterleaved)
\end{lstlisting}

\section{Результат}
\begin{verbatim}
Перемежение:
  После перемежения: 570 бит
  Первые 10 бит до перемежения: 1011000101
  Первые 10 бит после перемежения: 0010111010
\end{verbatim}
При обратном перемежении битовая последовательность полностью восстанавливается.

\section{Вывод}
Перемежитель и деперемежитель успешно реализованы. Псевдослучайная перестановка бит позволяет бороться с пакетными ошибками.